\section{实验设计、证据与结果记录}\label{sec:experiments}
\subsection{验证层级与当前事实}
本报告不将“已编写”“离线通过”“局部实机”和“完整验收”混为一谈。测试数量随仓库更新会变化，自动化测试通过也不能证明实际无碰撞。正式结果须绑定代码提交、设备、参数、场地与原始证据。
\begin{longtable}{p{32mm}p{54mm}p{51mm}}
\caption{已有证据及其范围}\\\toprule
项目 & 当前证据 & 能说明/不能说明\\\midrule\endfirsthead
\toprule 项目 & 当前证据 & 能说明/不能说明\\\midrule\endhead
动作副本 & 数据库帧、时间与非手臂通道核对；独立副本生成脚本 & 说明组合遵循指定规则；不证明抱持稳定性\\
人形局部路线 & 最新README记录抱起、走到A并约39.6\,cm处左转 & 支持首段局部运行；不证明B/C路线完成\\
人形离线验证 & 合成标记透视图与给定步态模型路线演练 & 说明设定条件下算法逻辑；不是真实误差分布\\
相机链 & 历史恢复检查连续取得15帧640\,$\times$\,480新图像 & 说明当次取流恢复；不保证持续运动不掉线\\
TurtleBot基础 & 小组确认底盘启动与键盘遥控可用 & 说明基础驱动可用；不等于自主导航成功\\
方式二隔离检查 & Kinetic构建、launch解析、缺少必要输入时零速度 & 部分软件与停止逻辑；未完成实地绕障验收\\
现场视频 & 新增235.52\,s原片及带时间戳截图，见下一节 & 支持部分运行过程；不覆盖完整放置与末车停稳\\\bottomrule
\end{longtable}

\subsection{分阶段实验与异常注入}
\begin{enumerate}
\item \textbf{单动作与视觉。}人工监督下分别测试站立、抓取、抱物前进、左右转和放置；固定相机姿态后测量标记距离/方位，统计多次误差与失效距离。
\item \textbf{人形单段与完整路线。}先测直行偏移，再测A左转/B右转，最后无人工干预完成A/B/C全程。分别记录动作次数、转角偏差、接触、掉落及相机断线。
\item \textbf{地图和单车导航。}确认35\,cm障碍进入原始深度与扫描，再检查地图。先静态目标导航、定位丢失和规划失败，再接入移动AR目标。
\item \textbf{目标与任务失效。}遮挡标签、停止目标消息、停止任务心跳、断开扫描或使TF过期，逐项测量停止延迟。人形最新程序的盲行与近距丢码启发式应单独评估。
\item \textbf{多车联调。}固定队列拓扑、共享地图与时间基准，避免速度控制源冲突；记录各车到达终点条件，不只记录领行机器人结束。
\end{enumerate}

\subsection{指标定义}
距离标定用均方根误差$\operatorname{RMSE}_d=\sqrt{N^{-1}\sum_i(\hat d_i-d_i)^2}$；$d_i$必须是与PnP输出同一定义的真实距离。跟随间距误差定义为$e_f(t)=d(t)-d_0$，报告平均绝对误差与最大偏差，同时说明$d$采用欧氏距离还是路径弧长。完整任务成功率定义为$N_{\rm success}/N_{\rm trials}$，每次试验均需记录，不能剔除失败试跑而不说明。

对接触次数、掉落次数、目标丢失持续时间、最长无新帧时间、停止延迟和各车到达时间分别统计。转角是否达到路线要求应通过外部测量或可复核视频判定，不能把“看见下一张码”直接记为精确90度。

\begin{table}[H]\centering\small
\caption{正式结果汇总模板；破折号表示未测，不表示零}
\begin{tabularx}{\linewidth}{Yp{26mm}p{26mm}p{25mm}}\toprule
项目 & 样本数 & 结果 & 证据编号\\\midrule
标记距离误差 & 待测 & --- & 待填写\\
满载单步距离/偏航 & 待测 & --- & 待填写\\
人形完整路线成功率 & 待测 & --- & 待填写\\
方式二完整跟随成功率 & 待测 & --- & 待填写\\
35\,cm障碍检测与避让 & 待测 & --- & 待填写\\
目标/任务失效停止延迟 & 待测 & --- & 待填写\\
全队到达率/总时长 & 待测 & --- & 待填写\\\bottomrule
\end{tabularx}\end{table}

此表用于统一后续试验的统计口径，当前空缺保持为未测状态。填入结果时应为每次试验绑定版本、参数、日志和视频，失败试次同样保留，以免只报告成功样本而高估系统性能。


\subsection{数据归档与照片使用}
运行记录模板保存在\code{latex/data/experiment_records.csv}。每行对应一次试验，保存代码版本、设备、运行参数、结果及证据路径；原始视频和bag较大时放在团队共享存储中，报告记录可访问的证据位置。不要将不同机器人的动作参数或不同版本的相机设置合并统计。

用户提供的通道、实机窗口及TurtleBot照片可作为场地和调试过程说明；应将原图归档到\code{figures/site}、\code{figures/humanoid}和\code{figures/follow2}并注明出处。本次已归档赛道原图及比赛视频原始帧，替换人形调试窗口和底盘实物图位；来源与时间戳见下一节。截图不代替连续运行记录。
